\chapter{Conventions at a glance}
\section*{Levels and indices}
\begin{tabular}{@{}ll@{}}\toprule
Input & Meaning\\\midrule
\fn{c[i]} & a spinful level, without operator-type or spin index\\
\fn{c[CR,i,UP]} & creation of its up electron\\
\fn{c[AN,i,DO]} & annihilation of its down electron\\
\fn{number[c[i]]} & sum over spin of the occupations\\
\fn{hubbard[c[i]]} & $n_{i\uparrow}n_{i\downarrow}$ without $U$\\
\fn{spinplus[c[i]]} & $c^\dagger_{i\uparrow}c_{i\downarrow}$\\
\fn{S[i,4]}, \fn{S[i,5]} & raising and lowering of a declared spin\\
\bottomrule\end{tabular}
\par\medskip
The internal type indices are \fn{CR}=0 and \fn{AN}=1; spin-$1/2$
labels are \fn{UP}=1, \fn{DO}=0. Do not infer bit order by sorting
these integers: consult the current \fn{BASIS}. A bare boson normally
needs no spin index. For spinless fermions use the mode convention of
the constructor in question; the generated explicit operator may
carry a trailing zero.

\section*{Order, charge and conjugation}
Normal ordering is a rule for reducing an operator, not an expectation
value. \fn{conj} reverses the order of factors and changes their type.
Complex coefficients must be conjugated too. The empty state has
$Q=-L$ for $L$ spin-degenerate levels; half filling has $Q=0$.
Highest-weight spin bases represent each complete spin multiplet once.
Reflection parity acts on states including the fermionic permutation
sign. Fermion parity $(-1)^N$ and spatial-reflection parity need not
agree. Energy gaps must specify which quantum numbers are held fixed.

\section*{Bloch, Nambu and retarded conventions}
An annihilator in real space has Fourier phase $+ikR$ and normalization
$N_s^{-1/2}$. Ordered orbitals within a cell fix the rows of $h(k)$.
Including their positions changes the phases of off-diagonal elements,
but neither the eigenvalues nor gauge-invariant observables. In the
particle-first Nambu spinor the hole block uses $-h^T(-k)$ for fermions;
the prefactor is $1/2$ when summing the full doubled basis.

For $\operatorname{Im}z>0$ the transform $e^{izt}$ of the retarded
kernel gives poles at $E_n-E_m$, with denominator $z-(E_n-E_m)$.
Fermionic diagonal spectral density integrates to one per canonical
mode including both signs of frequency. Use $\mathcal L A=[A,H]$;
exchanging the order of the commutator reflects all pole frequencies.
